In general, we will follow standard mathematical conventions with respect to notation.
For example, the union of sets is written as $\cup$, the intersection as $\cap$, and the element and subset relations as $\in$ and $\subset$, respectively.
However, for the sake of our convenience and hopefully better readability, we will fix some symbols for variables in the table at the end of this section.
Furthermore, we will adapt the notation given in the definitions.
\\
Our first definition is about a certain set $\mathbf{C}$, which always has a coordinate $\mathrm{ob}_{\mathbf{C}}$.
What we mean by adapting this notation is, for example, that we always write $\mathrm{ob}_{\mathbf{C}}$ for this coordinate - regardless of whether $\mathbf{C}$ is another token (in which case the index of $\mathrm{ob}_{\mathbf{C}}$ is naturally replaced by the new token).
If an explicit set is given a name, it will retain its name.
For example, the category of sets will always be written as $\mathbf{Set}$, and vice versa, this token always stands for the category of sets.
\\
Phrases like {\glqq}we (usually) write{\grqq} also indicate that the notation is fixed throughout the text.
Overall, we've tried to choose a consistent notation in the sense of using {\glqq}different symbol styles for different things{\grqq}.
We also try to use notation that doesn't leave much room for ambiguity.
But this can get quite cumbersome.
Therefore, we usually allow for more concise notations using the $\doteq$ token, even if the notation then becomes a bit inconsistent.
The strictly correct notation is on the right of the $\doteq$, while the more convenient (but perhaps inconsistent) one is on the left.
As an example of what we mean, suppose we have a thing depending on some set $Y$, which we denote as $T_{Y}$.
If the set is clear from context, it makes sense to just write $T$ for the thing.
We indicate that this is allowed by writing
\begin{align*}
  T
  &\doteq
  T_{Y}
\end{align*}
Last but not least, note: if undefined tokens like $f_{12}$ appear in the text, we strongly recommend referring to the following table first.
\begin{longtable}{cl}
  \hline
  Symbol
  &
  Meaning
  \\
  \hline
  \endhead
  \hline
  \endfoot
  $\mathbb{N}_{n}$
  &
  set of natural numbers $\lbrace 0,\ldots,n \rbrace$
  \\[5pt]
  $\mathbb{N}_{n}^{\times}$
  &
  set of natural numbers $\lbrace 1,\ldots,n \rbrace$
  \\[5pt]
  $\mathbb{N}_{0}^{\times}$
  &
  other notation for empty set $\emptyset$
  \\[5pt]
  $\mathbb{G}$
  &
  set of lowercase Greek letters $\alpha,\beta,\ldots,\omega$
  \\[5pt]
  $\mathcal{U}$
  &
  a variable for a Grothendieck universe
  \\[5pt]
  $K$
  &
  variable for a small set
  \\[5pt]
  $K_{a}$
  &
  variable for a small set for all $a \in \mathbb{G}$
  \\[5pt]
  $Y$
  &
  variable for a set
  \\[5pt]
  $\mathfrak{P}(Y)$
  &
  power set of $Y$
  \\[5pt]
  $\mathrm{pr}_{n}^{N}$
  &
  projection to the $n$-th coordinate of an $N$-tuple
  for $n \leq N$ and $n,N \in \mathbb{N}^{\times}$
  \\[5pt]
  $\mathbf{C}$
  &
  variable for a category
  \\[5pt]
  $X$
  &
  variable for an object of $\mathbf{C}$
  \\[5pt]
  $X_{n}$
  &
  variables for objects of $\mathbf{C}$ for all $n \in \mathbb{N}_{9}$
  \\[5pt]
  $f_{n_{1}n_{2}}$
  &
  variables for morphisms from $X_{n_{1}}$ to $X_{n_{2}}$ in
  $\mathbf{C}$ for all $n_{1},n_{2} \in \mathbb{N}_{9}$
  \\[5pt]
  $\mathbf{C}_{a}$
  &
  variables for categories for all $a \in \mathbb{G}$
  \\[5pt]
  $X^{a}$
  &
  variable for an object of $\mathbf{C}_{a}$ for all $a \in \mathbb{G}$
  \\[5pt]
  $X_{n}^{a}$
  &
  variables for objects of $\mathbf{C}_{a}$
  for all $n \in \mathbb{N}_{9}$ and all $a \in \mathbb{G}$
  \\[5pt]
  $f_{n_{1}n_{2}}^{a}$
  &
  variables for morphisms from $X_{n_{1}}^{a}$ to $X_{n_{2}}^{a}$
  in $\mathbf{C}_{a}$
  for all $n_{1},n_{2} \in \mathbb{N}_{9}$ and all $a \in \mathbb{G}$
  \\[5pt]
  $\mathbf{J}$
  &
  variable for a category
  \\[5pt]
  $J$
  &
  variable for an object of $\mathbf{J}$
  \\[5pt]
  $J_{n}$
  &
  variables for objects of $\mathbf{J}$ for all $n \in \mathbb{N}_{9}$
  \\[5pt]
  $j_{n_{1}n_{2}}$
  &
  variables for morphisms from $J_{n_{1}}$ to $J_{n_{2}}$ in
  $\mathbf{J}$ for all $n_{1},n_{2} \in \mathbb{N}_{9}$
  \\[5pt]
  $F_{a_{1}a_{2}}$
  &
  variable for a functor from $\mathbf{C}_{a_{1}}$ to $\mathbf{C}_{a_{2}}$
  for all $a_{1},a_{2} \in \mathbb{G}$
  \\[5pt]
  $F_{n}$
  &
  variables for functors from $\mathbf{C}$ to $\mathbf{C}_{\alpha}$
  for all $n \in \mathbb{N}_{9}$
  \\[5pt]
  $\mathsf{T}_{n_{1}n_{2}}$
  &
  variables for natural transformations from $F_{n_{1}}$ to $F_{n_{2}}$
  for all $n_{1},n_{2} \in \mathbb{N}_{9}$
  \\[5pt]
  $\mathcal{M}_{\mathbf{C}}$
  &
  variable for a monoidal category
  $(\mathbf{C},\otimes,\mathsf{A},1,\mathsf{L},\mathsf{R})$ 
  \\[5pt]
  $\mathcal{M}_{\mathbf{C}_{\alpha}}$
  &
  variables for monoidal categories
  $(\mathbf{C}_{a},\otimes_{a},\mathsf{A}_{a},1_{a},
  \mathsf{L}_{a},\mathsf{R}_{a})$
  for all $a \in \mathbb{G}$
  \\[5pt]
  $\mathcal{B}_{\mathbf{C}}$
  &
  variable for a braided monoidal category
  $(\mathbf{C},\otimes,\mathsf{A},1,\mathsf{L},\mathsf{R},\mathsf{B})$ 
  \\[5pt]
  $\mathcal{B}_{\mathbf{C}_{\alpha}}$
  &
  variables for braided monoidal categories
  $(\mathbf{C}_{a},\otimes_{a},\mathsf{A}_{a},1_{a},
  \mathsf{L}_{a},\mathsf{R}_{a},\mathsf{B}_{a})$
  for all $a \in \mathbb{G}$
  \\[5pt]
  $\mathcal{M}_{F_{a_{1}a_{2}}}$
  &
  variables for monoidal functors
  $(F_{a_{1}a_{2}},\mathsf{H}_{a_{1}a_{2}},\Phi_{a_{1}a_{2}})$
  \\
  $\phantom{\mathcal{M}_{F_{a_{1}a_{2}}}}$
  &
  from $\mathcal{M}_{\mathbf{C}_{a_{1}}}$ to $\mathcal{M}_{\mathbf{C}_{a_{2}}}$
  for all $a_{1},a_{2} \in \mathbb{G}$
  \\[5pt]
  $\mathcal{B}_{F_{a_{1}a_{2}}}$
  &
  variables for braided monoidal functors
  $(F_{a_{1}a_{2}},\mathsf{H}_{a_{1}a_{2}},\Phi_{a_{1}a_{2}})$
  \\
  $\phantom{\mathcal{B}_{F_{a_{1}a_{2}}}}$
  &
  from $\mathcal{M}_{\mathbf{C}_{a_{1}}}$ to $\mathcal{M}_{\mathbf{C}_{a_{2}}}$
  for all $a_{1},a_{2} \in \mathbb{G}$
  \\[5pt]
  $\mathcal{M}_{F_{n}}$
  &
  variables for monoidal functors $(F_{n},\mathsf{H}_{n},\Phi_{n})$
  from $\mathbf{C}$ to $\mathbf{C}_{\alpha}$ for all $n \in \mathbb{N}_{9}$
  \\[5pt]
  $\mathcal{B}_{F_{n}}$
  &
  variables for braided monoidal functors $(F_{n},\mathsf{H}_{n},\Phi_{n})$
  from $\mathbf{C}$ to $\mathbf{C}_{\alpha}$ for all $n \in \mathbb{N}_{9}$
  \\
  \hline
\end{longtable}
